%!TEX root = ./../../springer-mono-format/book.tex
%!TEX root = ./../../thesis-trier.tex

Here, we present new characterizations
of primitive permutation groups by properties
related to synchronization and subset mapping properties of associated
automata having a defect-one letter and containing the group in their
transformation semigroups.
Specifically, \cref{sync_max:thm:prim}
characterizes primitivity
by complete reachability, complete distinguishability,
sync-maximality, and distinguishability of pairs of disjoint $k$-sets for some $4 \le k^2 < n$ in associated automata with $n$ states.
\Cref{sync_max:cor:prim_grp}
formulates the results from the present section
in a purely algebraic language, without
reference to automata. The results of the present section
are taken from~\cite{HoffmannLATAextended2}. 
The following lemma is used in the proof of \cref{sync_max:thm:prim}.
 
\begin{lemma} \label{lem:fivetrans} \label{lem:threetrans}
	Let $G \le_G \mathcal S_n$.	Then~$G$ is transitive if one of the following holds:
	
	\begin{enumerate}
	\item \label{lem:fivetrans:1}  We have $n \ge 5$ and for each 
	simple idempotent $t \colon [n] \to [n]$ every $k$-set of states
	is distinguishable from every disjoint $l$-set of states
	$\langle G \cup \{t\}\rangle_S$
	for some $k,l \ge 2$ with $k + l \le n$.\footnote{In~\cite{HoffmannLATAextended2} this was only proved for $k = l = 2$.}
	
	\item \label{lem:fivetrans:2}  We have $n \ge 3$
	and for each 
	simple idempotent $t \colon [n] \to [n]$  
	every $k$-set of states is distinguishable from every distinct $l$-set of states
	in $\langle G \cup \{t\}\rangle_S$
	for some $2 \le k, l \le n$ where $n > k$ or $n > l$.
	\end{enumerate}
\end{lemma}
\begin{proof}
\ref{lem:fivetrans:1}: Suppose~$G$ is not transitive. By \cref{syncmax:prop:intransitive-subset} there exists a subset $A \subseteq [n]$
 permuted by each element from~$G$ containing three distinct $a_1, a_2, a_3 \in A$
 and such that there exists $b \in [n] \setminus A$. 
 First, suppose the complement of~$A$ contains two distinct points.
 Consider the simple idempotent $t \colon [n] \to [n]$
 mapping one point from $[n] \setminus A$
 to another point in $[n] \setminus A$.
 Then we have $h(a_1) \ne h(a_2)$
 and $h(a_3) \ne h(b)$ for all $h\in \langle G \cup \{t\} \rangle_S$.
 Since $k + l \le n$, we can pick $A, B \subseteq [n]$ where $|A| = l$ and $|B| = k$ such that $A \cap B = \emptyset$
 with $a_1, a_2 \in A$ and $a_3, b \in A$.
 Then~$A$ and~$B$ are not distinguishable.\footnote{If the complement of~$A$ contains two distinct points $b_1, b_2$, the same conclusion can be drawn by considering $\{a_1, b_1\}$ and $\{a_2, b_2\}$ instead of $\{a_1, a_2\}$ and $\{a_3, b\}$.}
 Next, suppose $[n] \setminus A = \{b\}$.
 Since~$n \ge 5$, there exists a fourth point $a_4 \in A \setminus\{a_1, a_2, a_3\}$.
 Consider the simple idempotent $t \colon [n] \to [n]$
 mapping $b$ to~$a_1$.
 We have $h(a_1) \ne h(a_2)$
 and $h(a_3) \ne h(a_4)$ for all $h\in \langle G \cup \{t\} \rangle_S$
 since $t(A) = A$ and $g(A) = A$ for $g \in G$.
 Pick $A, B \subseteq [n]$ where $|A| = l$ and $|B| = k$ such that $A \cap B = \emptyset$
 with $a_1, a_2 \in A$ and $a_3, a_4 \in A$.
 Then~$A$ and~$B$ are not distinguishable.
 
\ref{lem:fivetrans:1}: Suppose~$G$ is not transitive. By \cref{syncmax:prop:intransitive-subset} 
  there exists a subset $A \subseteq [n]$
  permuted by each element from~$G$
  and three distinct $a,b,c \in [n]$ with $a,b \in A$ and $c \in [n] \setminus A$.
  Let $t \colon [n] \to [n]$ be the simple idempotent mapping~$a$ to~$b$.
  We have $h(a) \ne h(c)$
  and $h(b) \ne h(c)$ for all $h \in \langle G \cup \{t\} \rangle_S$
  since $t(A) \subseteq A$, $t([n] \setminus A) = [n] \setminus A$, $g(A) = A$
  and $g([n] \setminus A) = [n] \setminus A$ for $g \in G$.
  Since $2 \le k,l \le n$ with $n > l$ or $n > k$ there 
  exist two distinct $A,B \subseteq [n]$ such that $|A| = k$ and $|B| = l$
  with $a,c \in A$ and $b,c \in B$,
  and these two subsets are indistinguishable
  since neither can be mapped to a singleton.\footnote{If $k > 2$
  	or $l > 2$, then we can also consider distinct $A, B \subseteq [n]$
  	with $a,c \in A \cap B$; however, the point~$b$
  	is needed for the case $k = l = 2$.}
\end{proof}

\begin{theorem}[Hoffmann~\cite{HoffmannLATAextended2}]
	\label{sync_max:thm:prim} 
	Let $G$ be a permutation group on the finite set~$Q$ of size at least five.
	Then the following are equivalent.
	\begin{enumerate}
		\item \label{sync_max:thm:prim:item:prim} $G$ is primitive. 
		
		\item  \label{sync_max:thm:prim:item:SI} Every SI-automaton $\semiautomat{A}{Q}{\Sigma}{\delta}$
		with a single simple idempotent letter, i.e.,~$|\Sigma_I| = 1$, and 
		$G = \{ \fword{\delta}{u} \mid u \in \Sigma_G^* \}$ is synchronizing. 
		
		\item  \label{sync_max:thm:prim:item:sync}  Every $\semiautomat{A}{Q}{\Sigma}{\delta}$ 
		having a letter of defect one with $G \subseteq \TS{\mathcal A}$ is synchronizing.
		
		\item  \label{sync_max:thm:prim:item:cr}  Every $\semiautomat{A}{Q}{\Sigma}{\delta}$ having
		a letter of defect one with $G \subseteq \TS{\mathcal A}$ is completely reachable.
		
		\item  \label{sync_max:thm:prim:item:cd}  Every $\semiautomat{A}{Q}{\Sigma}{\delta}$ having
		a letter of defect one
		with $G \subseteq \TS{\mathcal A}$ is completely distinguishable.
		
		\item  \label{sync_max:thm:prim:item:sm}  Every $\semiautomat{A}{Q}{\Sigma}{\delta}$ having
		a letter of defect one
		with $G \subseteq \TS{\mathcal A}$ is sync-maximal.	
		
		\item  \label{sync_max:thm:prim:item:ksets} For $\semiautomat{A}{Q}{\Sigma}{\delta}$ having
		a letter of defect one
		with $G \subseteq \TS{\mathcal A}$ every $k$-sets of states is distinguishable from every distinct $l$-set of states for some $2 \le k, l \le|Q|$ with $|Q| > k$ or $|Q| > l$.
			
		\item  \label{sync_max:thm:prim:item:klsets} For $\semiautomat{A}{Q}{\Sigma}{\delta}$ having
		a letter of defect one
		with $G \subseteq \TS{\mathcal A}$ every $k$-sets of states is distinguishable from every disjoint $l$-set of states for some $k, l \ge 2$
		with $kl < |Q|$.
	\end{enumerate} 
\end{theorem} 
\begin{proof}  	
The equivalences of Items~\ref{sync_max:thm:prim:item:prim}, \ref{sync_max:thm:prim:item:SI}, and~\ref{sync_max:thm:prim:item:sync}
are basically a reformulation of Rystsov's theorem (\cref{sync_max:thm:prim_rystsov})
in automata-theoretical language (see \cref{prelim:remarks_trans_semigrp_semiautomata}
for a discussion of the relation between transformation semigroups and semiautomata).

Next, if Items~\ref{sync_max:thm:prim:item:cr}, \ref{sync_max:thm:prim:item:cd}, or~\ref{sync_max:thm:prim:item:sm} hold,
then the automaton~$\mathcal A$ is synchronizing
by \cref{prelim:lem:compl_reach_implies_SCC_sync},\cref{sync_max:prop:dist_sync}, \crefitem{sync_max:thm:char_sync_max_aut}{sync_max:thm:char_sync_max_aut:1},
If \cref{sync_max:thm:prim:item:klsets,sync_max:thm:prim:item:ksets} holds,
then, by \cref{lem:fivetrans}, $G$ is transitive.
By restricting the alphabet, we can assume~$\mathcal A$ has a single non-permutational letter
of defect one. Then, with \cref{cor:defectatmostone-weaker} 
we get that~$\mathcal A$ is synchronizing. Thus, each of these assumptions
implies \cref{sync_max:thm:prim:item:sync}.\footnote{We can also give 
	an alternative argument that shows \cref{sync_max:thm:prim:item:prim}.
From assuming \cref{sync_max:thm:prim:item:ksets}: 
Suppose there exists a non-trivial $G$-congruence $\sim$.
Then there exist distinct $a,b,c \in X$
such that $a \sim b$
and $a \not\sim c$.
Consider the simple idempotent $t \colon X \to X$
mapping $a$ to $b$.
Then $\sim$ is a $\langle G \cup \{t\} \rangle_S$-congruence
and so $h(a) \not\sim h(c)$
and $h(b) \not\sim h(c)$ for all $h \in \langle G \cup \{t\} \rangle_S$.
Thus, no two subsets $A,B \subseteq X$
with $\{a,c\} \subseteq A$ and $\{b,c\} \subseteq B$
are distinguishable. 

A similar argument works when assuming \cref{sync_max:thm:prim:item:klsets}.
If we can pick distinct $a,b,c,d$
with $a \sim b$, $c \not\sim a$, and $d \not\sim a$.
then we can conclude as before that disjoint sets $A,B$
with $\{a,b\}\subseteq A$ and $\{c,d\} \subseteq B$ are not distinguishable.
Otherwise, if four distinct $a,b,c,d$ as stated cannot be picked, then $G$ has a fixpoint.
In this case, consider the simple idempotent that maps this fixpoint
to some other point. Then pairs of subsets where neither contains
the fixpoint are not distinguishable since it is not possible
to map them to smaller sets.}

Lastly, suppose~$G$ is primitive.
Let $\semiautomat{A}{Q}{\Sigma}{\delta}$ 
with $G \subseteq \TS{\mathcal A}$ having a letter of defect one.
By \cref{sync_max:prop:prim_implies_transitive}, $G$ is transitive.
By \cref{sync_words:lem:idempotent} we find that there exists a word~$u$
inducing a simple idempotent. Now consider a new automaton on the same state set
but where $u$ and all permutational letters are regarded as letters.
This gives an SI-automaton that also has~$G$ in the group of units of its transformation monoid.
Since we have already established the equivalence of Items~\ref{sync_max:thm:prim:item:prim} 
and~\ref{sync_max:thm:prim:item:sync}, we find that this SI-automaton is synchronizing.
By \cref{syncmax:cor:SI-transitive-grp-of-units} it is also completely reachable
and completely distinguishable. The transformation monoid of this SI-automaton is a submonoid of $\TM{\mathcal A}$. 
This yields that~$\mathcal A$
is completely reachable and completely distinguishable.
Thus we obtain
Items~\ref{sync_max:thm:prim:item:cr},\ref{sync_max:thm:prim:item:cd},\ref{sync_max:thm:prim:item:sm},
\ref{sync_max:thm:prim:item:ksets}, and~\ref{sync_max:thm:prim:item:klsets}.
This concludes the proof and establishes all the equivalences.
\end{proof}

\begin{corollary} \label{syncmax:cor:prim}
	Let $\semiautomat{A}{Q}{\Sigma}{\delta}$
	with at least five states 
	such that the group of units of $\TS{\mathcal A}$ is primitive.
	Then $\mathcal A$ is completely reachable if and only if it is sync-maximal.	
\end{corollary}
\begin{proof}
	If $\mathcal A$ is completely reachable, it must contain a letter $a$ of defect one.
	By \cref{sync_max:thm:prim},
	the automaton restricted to the alphabet
	$\{a\} \cup \{ x \in \Sigma \mid \text{$x$ is a permutational letter} \}$
	is completely distinguishable.
	Thus, $\mathcal A$
	is completely distinguishable
	and so sync-maximal.	
	Conversely, since a primitive group is transitive by \cref{sync_max:prop:prim_implies_transitive}, 
	$\mathcal A$ is strongly connected and so sync-maximality
	implies complete reachability.		
\end{proof}

By \cref{syncmax:cor:prim}, if the transformation monoid contains a primitive permutation group,
complete reachability implies complete distinguishability.
But the converse implication does not hold in general,
see, for example, \cref{sync_max:ex:compl-dist-non-cr}.

The equivalences from \cref{sync_max:thm:prim}
also hold true when restricted to the following
subclasses of automata: SI-automata with $\Sigma_I \ne \emptyset$, 
SI-automata with $|\Sigma_I| = 1$, automata with a group of units equal to $G$
having a letter of defect one, and automata having only a single 
non-permutational letter of defect one.


\label{def:Afa}
For $\semiautomat{A}{Q}{\Sigma}{\delta}$,  $t \colon Q \to Q$,
and a new symbol $a \notin \Sigma$
we define a semiautomaton $\semiautomat{A_{t,a}}{Q}{\Sigma \cup \{a\}}{\mu}$
with, for $q \in Q$ and $x \in \Sigma \cup \{a\}$, \[
\applytrans{\mu}{q}{x} = \begin{cases} 
	\applytrans{\delta}{q}{x} & \text{if } x \in \Sigma, \\ 
	t(q) & \text{if } x = a.
\end{cases}
\]
\cref{sync_max:fig:equivalences_perm_aut}
shows the main equivalences from \cref{sync_max:thm:prim_min,sync_max:thm:prim}
when adjoining a single letter of defect one.



\begin{figure}[ht]
	\centering
	
	\begin{tikzpicture}[
		node distance=3.5em and 4em,
		outer sep=4pt,
		every node/.style={
			draw,
			rectangle,
			align=center,
			minimum width=16em,
			minimum height=3em
		}
		]
		\node (A) [draw, rectangle, align=center]
		{$\TS{\mathcal A}$ contains primitive\\ permutation group};
		
		\node (A2) [right=of A]
		{$\forall t \colon Q \to Q \text{ of defect one,}\;
			\forall a\notin\Sigma:$\\
			$\mathcal A_{t,a}\text{ completely reachable.}$};
		
		\node (B) [draw, rectangle, align=center, below=of A]
		{$\forall q\in Q\;
			\forall F\in\mathcal P(Q)\setminus\{\emptyset,Q\}:$\\
			$\mathcal A_{q,F}\text{ is minimal.}$};
		
		\node (C) [draw, rectangle, align=center, below=of A2]
		{$\forall t\colon Q\to Q\text{ of defect one,}\;
			\forall a\notin\Sigma:$\\
			$\mathcal A_{t,a}\text{ sync-maximal.}$};
		
		\draw[<->, thick] (A) -- (B);
		\draw[<->, thick] (A) -- (C);
		\draw[<->, thick] (B) -- (C);
		\draw[<->, thick] (A) -- (A2);
		\draw[<->, thick] (A2) -- (C);
		\draw[<->, thick] (A2) -- (B);
	\end{tikzpicture}
	\caption[Equivalences for automata with primitive group in transformation semigroup]{Equivalences  from \cref{sync_max:thm:prim_min,sync_max:thm:prim} for $\semiautomat{A}{Q}{\Sigma}{\delta}$.}
	\label{sync_max:fig:equivalences_perm_aut}
\end{figure}
Next, as an example,
we give a formulation of some properties from \cref{sync_max:thm:prim}
without any reference to automata. In the sections
to follow all results will be given without reference
to automata, but they all can be formulated
as statements about automata similarly as in \cref{sync_max:thm:prim}. 
We combine two statements (related via \cref{sync_words:lem:idempotent})
into one by putting the word ``idempotent``
in square brackets, see \cref{syncmax:rem:square-brackets}.

\begin{corollary}[Hoffmann~\cite{HoffmannLATAextended2}]
\label{sync_max:cor:prim_grp}
	Let $G \le_G \mathcal S_n$ be a permutation group with $n \ge 5$. Then the following are equivalent.	
\begin{enumerate}
\item \label{sync_max:cor:prim:item:prim} $G$ is primitive.

\item \label{sync_max:cor:prim:item:compl_reach} 
 For every [idempotent] transformation $t \colon [n] \to [n]$ of defect one, $\langle G \cup \{t\} \rangle_S$ is completely reachable.

\item For every [idempotent] transformation $t \colon [n] \to [n]$ of defect one, $\langle G \cup \{t\} \rangle_S$ is completely distinguishable.

\item \label{sync_max:cor:prim_grp:4} For every [idempotent] transformation $t \colon [n] \to [n]$ of defect one, in $\langle G \cup \{t\} \rangle_S$ every $k$-set is distinguishable from every distinct $l$-set
for some $2 \le k,l < n$.

\item \label{sync_max:cor:prim_grp:5} For every [idempotent] transformation $t \colon [n] \to [n]$ of defect one, 
in $\langle G \cup \{t\} \rangle_S$ every $k$-set is distinguishable from every disjoint $l$-set
for some $k,l \ge 2$ with $kl < n$.
\end{enumerate}
\end{corollary}

We cannot omit the assumption $kl< n$ in Item~\ref{sync_max:cor:prim_grp:5} 
from \cref{sync_max:cor:prim_grp} as demonstrated by the following example with $k = l = 2$.

\todo[inline]{hier 4 in IC paper?}
\begin{example}
	Let $g = (0 ~ 1 ~ 2)(3)$. 
	Then~$G = \langle \{g\} \rangle_G$ is imprimitive as it is not transitive.	
    Set
	\begin{align*}
		S = \{ \{0,3\}^h \mid h \in G \} & = \{ \{0,3\}, \{1,3\}, \{2,3\} \}, \\
		T = \{ \{0,1\}^h \mid h \in G \} & = \{ \{0,1\}, \{1,2\}, \{2,0\} \}.
	\end{align*}
    Let $\{a,b\}$, $\{c,d\}$ be two disjoint $2$-sets.
    Then $\{a,b\} \in S$ and $\{c,d\} \in T$, or
    $\{a,b\} \in T$ and $\{c,d\} \in S$.
    Now, for $t \colon [4] \to [4]$ of defect one,
    there exists $h \in G$
    mapping precisely one of $\{a,b\}$ or $\{c,d\}$
    to the unique $2$-set
    mapped by $t$ to a singleton.
    Then $th$ distinguishes $\{a,b\}$ and $\{c,d\}$.
\end{example}